\documentclass[12pt,a4paper]{article}
\usepackage[utf8]{inputenc}
\usepackage[T1]{fontenc}
%\usepackage[french]{babel} % 若需法语支持
\usepackage[a4paper, top=3cm, bottom=3cm, left=2.5cm, right=2.5cm]{geometry}
\usepackage{amsmath,amssymb}
\usepackage{graphicx}
\usepackage{authblk} % 添加作者机构支持
\usepackage{hyperref}
\hypersetup{
    colorlinks=true,
    linkcolor=blue,
    citecolor=red,
    urlcolor=cyan
}
\bibliographystyle{elsarticle-num} % 参考文献样式
\usepackage{doi}
\usepackage{float}
\usepackage{caption}
\usepackage{lineno}
\usepackage{soul}
\usepackage{xcolor}
\usepackage{multirow}

% 自定义命令与符号
\def\dsp{\displaystyle}
\newtheorem{remark}{Remark}
\def\O{\Omega}
\def\Ff{{\cal F}}
\def\d{{\rm d}}
\def\red{\textcolor{red}}
\def\blue{\textcolor{blue}}
\def\kb{\textcolor{magenta}}
\def\jm{\textcolor{cyan}}



\begin{document}

\begin{titlepage}
    \centering
    \includegraphics[width=1\textwidth]{images/logo.jpg}\par\vspace{1.5cm}
    
    

    {\huge\bfseries M2 Internship Report\\Advanced Numerical Models in Hydrogeology\par}
    \vspace{1cm}
    
    {\large Internship report completed from May 13, 2025 to August 24, 2025\par}
    \vspace{0.5cm}

    \includegraphics[width=0.82\textwidth]{images/title.png}\par
    \vspace{1cm}

    \textbf{Supervisors:}
    Konstantin Brenner,
    Jean-Raynald de Dreuzy and
    Jean Marçais
    \vspace{0.5cm}
    
    \textbf{Authors:}
    Qinyan YANG



\end{titlepage}


\newpage

\vspace{1cm}
\tableofcontents

\newpage



\section{Introduction}  
In the context of climate change and evolving land use, accurately modeling the exchanges between the unsaturated and saturated zones is essential for water resource management and scientific research. The three-dimensional Richards equation provides a detailed description of these interactions, but its strong nonlinearity and high dimensionality make numerical solutions computationally demanding. In contrast, the two-dimensional Boussinesq equation—derived under the Dupuit assumption in the saturated zone—offers high numerical efficiency but does not account for unsaturated zone dynamics.\\  

To balance physical accuracy and computational performance, recent studies have proposed a coupled approach: applying the one-dimensional Richards equation to the unsaturated zone and the two-dimensional Boussinesq equation to the saturated zone. A preliminary implementation of this Boussinesq–Richards coupled model had already been developed by the supervising team. During this internship, my work was structured around three main tasks:  

\begin{enumerate}  
    \item \textbf{Model testing and bug reporting}  
    \vspace{0.25cm}\\
    I designed a series of representative test cases (steady-state and transient flows, varying soil characteristics, etc.) to evaluate the model’s numerical stability and robustness. During execution, I identified various issues (convergence failures, timing mismatches in recharge and outflow, etc.) and provided detailed reports—including error messages, reproducible procedures, and input files—to facilitate their resolution.  
    
    \item \textbf{Result validation and performance evaluation}\vspace{0.25cm}\\
    Validation commenced under steady-state, constant-recharge conditions where analytical solutions exist. The Boussinesq-Richards (BR) model's water table solutions were benchmarked against these closed-form results. A series of Python scripts was developed to: (1) extract saturation profiles, (2) draw water table profiles, and (3) quantify outflow rates. These outputs were systematically benchmarked against the 2D Richards (R) model, with deviations between the models calculated to verify the accuracy of the Boussinesq-Richards formulation across diverse scenarios.
    

    \item \textbf{Participation in scientific article preparation}\vspace{0.25cm}\\
    I contributed to drafting the manuscript describing the testing and validation of the Richards–Boussinesq model, including the preparation of figures and simulation results for publication. 
\end{enumerate}  

Through these contributions, this internship enhanced the reliability and usability of the Richards–Boussinesq coupled model and established a robust testing and validation protocol for its future application to large-scale hydrogeological problems.  


\section{Mathematical background}  

\subsection{Richards equation}  
Richards' equation serves as a \textbf{reference model} in groundwater hydrology because it can simultaneously describe flow both in the unsaturated zone and below the water table. Denoting by \( s \) and \( p \) the unknown fields of water saturation and pressure, the equation is expressed as:
\begin{equation}\label{richards}
\phi \partial_t s - \mathrm{div} \left( k(s) \nabla ( p + z ) \right) = f(p,z)
\end{equation}
where $s$ is the relative saturation, $p$ is the pressure, $z$ the vertical coordinate, and $\phi$ the soil porosity.  
The source term \( f(p,z) \) models the impact of precipitation, evapotranspiration, and other sources or sinks, while the relative conductivity \( k(s) \) is typically given by the law \( k(s) = K_s\,s^\alpha \), with \( \alpha > 1 \) and $K_s$ denoting the saturated hydraulic conductivity.  \\

The equation \eqref{richards} is complemented by an algebraic closure law \( s = S(p) \), where \( S \) is a continuous, nondecreasing function from \( \mathbb{R} \) to \( (0,1] \), describing the capillary properties of the medium \cite{van1980closed}.  
Physically, Richards' model is quite complete: it captures unsaturated–saturated transitions, capillary effects, and recharge dynamics. However, this generality comes at a high computational cost, especially in higher dimensions.

\subsection{Boussinesq equation}  
The Boussinesq (Dupuit–Forchheimer) model \cite{Dupuit,Forchheimer} simplifies Richards' equation by assuming hydrostatic pressure distribution in the saturated zone and neglecting capillary effects. Through vertical integration, this yields:
\begin{equation}\label{dupuit}
\partial_t (\gamma h) - \mathrm{div}_{xy} \left( K_s h \nabla_{xy}(h + z_b) \right) = r
\end{equation}
where \( h \) is the water table depth, \( z_b \) the bedrock elevation, \( \gamma \) the specific yield, and \( r \) the net recharge. The operators \( \mathrm{div}_{xy} \) and \( \nabla_{xy} \) denote horizontal divergence and gradient.  \\

Physically, this model is much lighter computationally and well adapted for large-scale aquifer dynamics. On the other hand, it cannot represent the unsaturated zone, infiltration fronts, or capillary storage processes, which limits its realism.  
Its derivation can be rigorously justified via asymptotic analysis of Richards’ equation under the assumption of vanishing aquifer depth or very high vertical conductivity \cite{Blendinger}.



\subsection{Coupled Boussinesq–Richards model}  
Our integrated approach combines both formulations through a dual-continuum representation, in which vertical columns are resolved with Richards’ equation and horizontally averaged dynamics are represented with a Boussinesq-type model.

\paragraph{Modified Richards equation:}  
\begin{equation}\label{mod_richards}
\phi \partial_t s - \partial_z \left( k(s) (\partial_z p + 1) \right) + \mathcal{F} = 0
\end{equation}
This equation is solved along full vertical columns, with the additional source term $\mathcal{F}$ representing the exchange with the regional water table.  
It thus preserves the ability of Richards’ equation to describe unsaturated flow and recharge dynamics.

\paragraph{Exchange flux:}  
\begin{equation}\label{exchange_flux}
\mathcal{F} = \frac{1}{\varepsilon} \left( p^+ - \overline{p}(h,z)^+ \right) \mathbf{1}_{\{z \leq h + z_b\}}
\end{equation}
where $\mathbf{1}$ denotes the indicator function restricting exchange to the saturated zone, and $\overline{p}(h,z)$ is the hydrostatic pressure induced by the water table.  
Physically, this formulation ensures a consistent coupling between local pressure dynamics in Richards’ columns and the regional water table represented by the Boussinesq model.

\paragraph{Modified Boussinesq equation (water table equilibrium):}  
\begin{equation}\label{mod_dupuit}
-\partial_x\left( h \partial_x \left( h + z_b \right) \right) - \int_{z_b}^{z_t} \mathcal{F}\, \mathrm{d}z = 0
\end{equation}
The integral term couples the vertical fluxes $\mathcal{F}$ from the Richards columns with the horizontal dynamics of the water table.  
This representation allows lateral groundwater flow to be simulated efficiently, while still accounting for the vertical interactions between saturated and unsaturated zones.  
As a result, the model combines the physical realism of Richards’ equation with the computational efficiency of the Boussinesq approach.



\subsection{Numerical discretization}

The coupled Boussinesq--Richards model is discretized using a fully implicit finite volume scheme.  
\begin{itemize}
    \item \textbf{Time discretization:} a backward Euler scheme is employed to guarantee unconditional stability.  
    \item \textbf{Space discretization:} the vertical Richards equation and the horizontal Boussinesq equation are discretized separately. Upwinding is used in the computation of fluxes to properly handle nonlinearities and to ensure robustness of the scheme.  
\end{itemize}

To efficiently solve the highly nonlinear Richards equation, a \emph{variable switching} technique is introduced. This approach employs an auxiliary variable that allows for a unified treatment of both saturated and unsaturated regimes, avoiding numerical difficulties associated with phase transitions. \\ 

The most crucial aspect of the discretization lies in exploiting the decoupling of the equations along spatial dimensions. A \emph{nonlinear elimination} method is used to reduce the coupled system to a smaller algebraic system involving only the water table height \(h\). This significantly improves computational efficiency and enables the use of fast iterative solvers such as Newton's method.






\section{Model testing and debugging}  

To evaluate the Richards–Boussinesq coupled model, a series of representative test cases were implemented, including both steady-state and transient flow conditions, as well as scenarios with varying soil hydraulic properties. The main goal was to assess the numerical scheme’s stability and verify the consistency of outputs under different conditions.\\ 

During these tests, several issues were identified:  

\begin{itemize}  
    \item \textbf{Convergence failures:} Three main sources of convergence problems were observed:  
    \begin{enumerate}  
        \item When recharge was defined as a step function, the strong nonlinearity introduced in the source term caused the Newton solver to diverge. To address this, the discontinuous step input was replaced with a smooth sigmoid-based approximation:  

        \begin{equation}
        R(t) = \sum_{i=1}^{N} A_i \left[ \frac{1}{1 + \exp\!\left(-(t-t_i)/\varepsilon\right)} - \frac{1}{1 + \exp\!\left(-(t-t_{i+1})/\varepsilon\right)} \right]
        \end{equation}  

        where $A_i$ is the amplitude of the $i$-th step, $t_i$ and $t_{i+1}$ are the start and end times of the step, and $\varepsilon$ controls the smoothness of the transition. Smaller $\varepsilon$ values produce sharper transitions, approaching a step, while larger values ensure smoother variations that facilitate Newton convergence. This approach preserves the physical behavior of recharge while preventing solver divergence. Figure~\ref{fig:sig} illustrates the resulting smooth recharge profile.  

        \begin{figure}[H]
            \centering
            \includegraphics[width=0.5\linewidth]{images/sig.png}
            \caption{Sigmoid-based smooth approximation\\
            of a step-type recharge function.}
            \label{fig:sig}
        \end{figure}

        \item For soils with extremely low hydraulic conductivity, Newton iterations occasionally failed to converge. This limitation was reported for further consideration.  

        \item In some scenarios, convergence was only achieved by automatically reducing the time step size. Implementing an adaptive time-stepping procedure allowed the solver to converge reliably.  
    \end{enumerate}  

    \item \textbf{Timing mismatches:} Two sources of mismatch between recharge input and aquifer response were identified:  
    \begin{enumerate}  
        \item Linear interpolation in the initial implementation distorted step-type recharge inputs, creating the appearance of a delayed system response. The sigmoid-based approximation corrected this issue by providing a smooth yet accurate representation of the recharge function.  

        \item A secondary mismatch arose in the time integration, where recharge and outflow were evaluated inconsistently. Aligning the flux calculation with the new implementation eliminated the temporal shift between recharge and aquifer response.  
    \end{enumerate}  
\end{itemize}  


Through this process, the main sources of numerical instability were identified and practical solutions—sigmoid smoothing, adaptive time stepping, and consistent flux evaluation—were implemented, significantly improving the robustness of the Richards–Boussinesq code and establishing a systematic debugging protocol for future development.
 

  

\section{Validation and performance analysis}


    %%%%%%%%%%%%%%%%%%%%%%%%%%%%%%%%%%%%%%%%%%%%%%%%%
    %%     Test cases with different soil types    %% 
    %%%%%%%%%%%%%%%%%%%%%%%%%%%%%%%%%%%%%%%%%%%%%%%%%

Table \ref{tab:brooks_corey_parametrization} presents the soil hydraulic parameters for the Brooks-Corey formulation \cite{brooks1964hydraulic}, applied to the loam and sand textural classes used in both the 2D Richards and Boussinesq-Richards models. These parameter values are sourced from \cite{rawls1982} and are consistent with the Hydrus soil library \cite{simunek2008b}.\\
\begin{table}[h]
    \centering
    \begin{tabular}{c|c c c c c c}
         Textural class & $\theta_r$ & $\theta_s$ & $\alpha$ & $n$ & $K_s$ \\
Unity & [$L^{3}L^{-3}$] & [$L^{3}L^{-3}$] & [$cm^{-1}$] & [$-$] & [$cm.day^{-1}$]  \\
 \hline
         Sand & 0.020 & 0.417 & 0.138 & 0.592 & 504 \\
         Loam & 0.027 & 0.434 & 0.0897 & 0.220 & 31.7\\
    \end{tabular}
    \caption{Hydraulic parameters for Brooks-Corey formulation}
    \label{tab:brooks_corey_parametrization}
\end{table}

%%%%%%%%%%%%%%%%%%%%%%%%%%%%%%%%%%%%%%%%%%%%%%%%%
%%               Analytical-solution           %% 
%%%%%%%%%%%%%%%%%%%%%%%%%%%%%%%%%%%%%%%%%%%%%%%%%
\subsection{Validation of Boussinesq-Richards Model Using Analytical Solutions}

Under steady-state conditions, the Boussinesq-Richards model can be described by the following ordinary differential equation (ODE):\vspace{0.3cm}
\begin{equation}
\frac{d}{dx} \left( \text{K}_s h \frac{d}{dx}(h + z_b) \right) = -R
\end{equation}
\vspace{0.25cm}

where \( h \) is the water table depth, \( z_b \) the bedrock elevation, and \( R \) represents the recharge rate.\\

To evaluate the performance of the Boussinesq-Richards model under steady-state conditions, we conducted a numerical simulation within a two-dimensional domain consisting of sand/loam soil. The vertical domain spans from \( z_b = 0\,\text{m} \) (bottom boundary) to \( z_t = 7\,\text{m} \) (surface elevation). The surface at \( z = z_t \) represents the ground surface and is treated as a seepage face, while the bottom boundary at \( z = z_b \) is defined as a no-flow boundary.\\

Horizontally, the domain extends from \( x = -25\,\text{m} \) to \( x = 25\,\text{m} \). The boundary at \( x = 25\,\text{m} \) is specified as a no-flow boundary, whereas the boundary at \( x = -25\,\text{m} \) simulates a river and is modeled using a Dirichlet boundary condition with a constant water level of \( 2\,\text{m} \).\\

Under these assumptions, the analytical solution for the water table elevation \( h(x) \) is given by:
\begin{equation}
    h(x) = \sqrt{ \frac{2R}{k} \left( x_{\text{max}} x - \frac{1}{2}x^2 \right) + h_0^2 - \frac{2R}{k} \left( x_{\text{min}} x_{\text{max}} - \frac{1}{2} x_{\text{min}}^2 \right) }
\end{equation}
Where, \(R\) is the constant recharge rate, set to 3 mm/day; \(k\) denotes the soil hydraulic conductivity; \(h_0\) is the prescribed head at the river boundary at \(x = x_{\text{min}}\); and \(x_{\text{min}}\) and \(x_{\text{max}}\) define the horizontal extent of the domain. 


\begin{figure}[H]
\centering
\makebox[\textwidth]{
\includegraphics[width=1.2\textwidth,
height=0.35\textwidth]{images/aVSn.jpg}
}
\caption{Comparison between the Boussinesq-Richards numerical solution and the analytical solution.}
\label{fig:ana_VS_num}
\end{figure}

As shown in Figure \ref{fig:ana_VS_num}, the numerical results obtained from the Boussinesq-Richards model exhibit excellent agreement with the analytical solution under steady-state, confirming the model’s accuracy under steady-state conditions.\\

\subsection{Validation Against 2D Richards Model}
\vspace{0.3cm}
To validate the proposed Boussinesq-Richards coupled approach, we conduct numerical experiments comparing it against the standard 2D Richards model. Both solvers are executed on identical test cases, with discrepancies quantified using hydraulic head and saturation error metrics. 

\subsubsection{Case 1: synthetic test cases with different soil types: sand and loam.}
\vspace{0.25cm}

    %%%%%%%%%%%%%%%%%%%%%%%%%%%%%%%%%%%%%%%%%%%%%%%%%
    %%               Boundary conditions           %% 
    %%%%%%%%%%%%%%%%%%%%%%%%%%%%%%%%%%%%%%%%%%%%%%%%%


For the first test case, the simulation domain used in this study is shown in Figure~\ref{fig:dom}, it consists of parallelepiped with a length of 50 meters, a vertical height of 7 meters, and an inclination of \text{5\%}. It is filled with sandy or loam materials. As illustrated, the domain includes two no-flow boundaries (along the bottom and right edges), two open boundaries (along the left edge), and a top boundary subject to an imposed flux. A Dirichlet boundary with a height of 2 meters is applied along the left edge to represent a lateral river. At the start of the simulation (initial state; see Supplementary S1.3), a constant recharge rate of 3 mm/day is applied until the system reaches a steady-state (shown in Figure~\ref{fig:steady_comparison}). This steady state is then used as the initial condition for subsequent simulations, during which time-varying recharge is imposed to compare the performance of the Boussinesq--Richards model and the Richards model.

\begin{figure}[H]
    \centering
    \includegraphics[width=0.9\textwidth]{images/0.png}
    \caption{Simulation domain for case 1}
    \label{fig:dom}
\end{figure}


%%%%%%%%%%%%%%%%%%%%%%%%%%%%%%%%%%%%%%%%%%%%%%%%%
%%                Initial-state                %% 
%%%%%%%%%%%%%%%%%%%%%%%%%%%%%%%%%%%%%%%%%%%%%%%%%
The initial states for sand and loam continuum are shown in Figures~\ref{fig:initial_C}. A constant recharge of 3 mm/day is first applied until steady-state is reached (Figure~\ref{fig:steady_comparison}).

\begin{figure}[H]
    \centering
    \includegraphics[width=1\textwidth]{images/0_1.png}
    \caption{\centering
    Initial saturation conditions. \\
    Boussinesq-Richards model for (a) sand and (b) loam continuum; \\
    2D Richards model for (c) sand and (d) loam continuum.
    }
    \label{fig:initial_C}
\end{figure}



%%%%%%%%%%%%%%%%%%%%%%%%%%%%%%%%%%%%%%%%%%%%%%%%%
%%                Steady-state                 %% 
%%%%%%%%%%%%%%%%%%%%%%%%%%%%%%%%%%%%%%%%%%%%%%%%%

Once steady-state (Figure~\ref{fig:steady_comparison}) is achieved, a variable vertical recharge is introduced, as shown in Figure~\ref{fig:variable_recharge}.

\begin{figure}[H]
    \centering
    \captionsetup{justification=centering}
    \includegraphics[width=1\textwidth]{images/1.png}
    \caption{Steady-state saturation. \\
    Boussinesq-Richards model for (a) sand and (b) loam continuum; \\
    2D Richards model for (c) sand and (d) loam continuum.}
    \label{fig:steady_comparison}
\end{figure}



%%%%%%%%%%%%%%%%%%%%%%%%%%%%%%%%%%%%%%%%%%%%%%%%%
%%            Recharge illustration            %% 
%%%%%%%%%%%%%%%%%%%%%%%%%%%%%%%%%%%%%%%%%%%%%%%%%
The variable vertical recharge consists of a 365 day time series with sharp transitions between recharge events, ranging from no recharge (0 mm/day) to 6 mm/day, as shown in Figure~\ref{fig:variable_recharge}.

\begin{figure}[H]
    \centering
    \includegraphics[width=0.45\textwidth]{images/2.png}
    \caption{Time-varying vertical recharge applied over 365 days.}
    \label{fig:variable_recharge}
\end{figure}

Based on the above setup, saturation distributions, water table positions, and outflow rates were obtained from 365 day simulations and are analyzed in the following section.\\



    %%%%%%%%%%%%%%%%%%%%%%%%%%%%%%%%%%%%%%%%%%%%%%%%%
    %%              mid saturation                 %% 
    %%%%%%%%%%%%%%%%%%%%%%%%%%%%%%%%%%%%%%%%%%%%%%%%%
To validate the Boussinesq--Richards model, its outputs are compared with those of the 2D Richards model. The analysis first focuses on the saturation profile by examining saturation values at $x=0$ (shown in Figure~\ref{fig:mid_results}) to highlight the differences between the two models. Four time steps were selected to capture the transient dynamics: the initial time, the final time (during the no-infiltration period), and two points within the infiltration phases. Since the simulation included three repeated infiltration events, the system responses during these periods are nearly identical. Therefore, the analysis concentrates on the time points near the beginning and the end of a single infiltration event (the second infiltration phase).

\begin{figure}[H]
\centering
\makebox[\textwidth]{%
    \includegraphics[width=1.15\textwidth, height=1\textwidth]{images/3.png}
}
\caption{\centering
Saturation in the middle of the domain.\\
Panel (a) for sand and (b) for loam continuum.
}
\label{fig:mid_results}
\end{figure}

As shown in Figure~\ref{fig:mid_results}, the sand continuum, due to its higher hydraulic conductivity compared to loam, exhibits more rapid saturation responses to recharge and drainage. This faster dynamic behavior, relative to the loam continuum, is consistent with the hydrogeological properties of the two soils. Overall, the saturation distributions predicted by the two models are in close agreement, the main differences appearing near the water table region. This observation is further quantified by the saturation error field, obtained as the difference between the saturation fields of the two models, and is presented in Figure~\ref{fig:Serr}.


\begin{figure}[H]
    \centering
    \includegraphics[width=0.6\linewidth]{images/serr.png}
    \caption{Saturation error field between the Boussinesq--Richards model and the 2D Richards model.}
    \label{fig:Serr}
\end{figure}


%%%%%%%%%%%%%%%%%%%%%%%%%%%%%%%%%%%%%%%%%%%%%%%%%
%%              Water Table                    %%
%%%%%%%%%%%%%%%%%%%%%%%%%%%%%%%%%%%%%%%%%%%%%%%%%

The water table positions are extracted at the same time points as the saturation analysis. The water table is defined as the vertical coordinate where the pressure equals zero. Figure~\ref{fig:wt_sand} and Figure~\ref{fig:wt_loam} illustrate the water table evolution for the sand and loam continua, respectively.\\ 

As shown, the sand continuum exhibits faster water table fluctuations in response to recharge and drainage due to its higher hydraulic conductivity. However, its relatively lower water retention limits the overall amplitude of these fluctuations. In contrast, the loam continuum, being more capillary, maintains a generally higher water table and responds more gradually to the infiltration events, showing larger water table variations. This behavior is consistent with the hydrogeological properties of the two soils.


\begin{figure}[H]
\centering
\makebox[\textwidth]{%
    \includegraphics[width=1\textwidth]{images/wt_sand.png}
}
\caption{
Water table positions in the sand continuum at selected time steps.
}
\label{fig:wt_sand}
\end{figure}

\begin{figure}[H]
\centering
\makebox[\textwidth]{%
    \includegraphics[width=1\textwidth]{images/wt_loam.png}
}
\caption{
Water table positions in the loam continuum at selected time steps.
}
\label{fig:wt_loam}
\end{figure}



    %%%%%%%%%%%%%%%%%%%%%%%%%%%%%%%%%%%%%%%%%%%%%%%%%
    %%                 errors                      %% 
    %%%%%%%%%%%%%%%%%%%%%%%%%%%%%%%%%%%%%%%%%%%%%%%%%
To further quantify the discrepancies between the Boussinesq–Richards and the 2D Richards models, two error metrics are introduced: a normalized mean absolute error for the saturation field and a similar metric for the water table position. These measures enable a systematic evaluation of both spatial and temporal differences between the two models.\\

The first metric evaluates the relative error in the saturation field, averaged over all grid cells:  
\begin{equation}
\text{Serr} = \frac{1}{NM} \sum_{i=1}^{N} \sum_{j=1}^{M} 
\left( \frac{\left| S^{\text{BR}}_{i,j} - S^{\text{R}}_{i,j} \right|}{S^{\text{R}}_{i,j}} \right) \times 100\% 
\end{equation}
where \( N \) and \( M \) denote the numbers of grid cells in the horizontal and vertical directions, respectively.  
Here, \( S^{\text{BR}}_{i,j} \) and \( S^{\text{R}}_{i,j} \) represent the saturation values at grid cell \( (i,j) \) 
for the Boussinesq–Richards and the 2D Richards models, respectively. This metric provides a spatially averaged measure of the relative saturation error expressed in percentage.\\

The second metric quantifies the relative deviation in water table elevations along the horizontal direction:  
\begin{equation}
\text{WTerr} = \frac{1}{N} \sum_{i=1}^{N} 
\left( \frac{\left| WT^{\text{BR}}_i - WT^{\text{R}}_i \right|}{WT^{\text{R}}_i} \right) \times 100\% 
\end{equation}
where \( N \) is the number of horizontal grid points, and  
\( WT^{\text{BR}}_i \) and \( WT^{\text{R}}_i \) denote the water table elevations at the \( i \)-th horizontal location 
for the Boussinesq–Richards and the 2D Richards models, respectively. This error metric captures the horizontal variability in water table predictions.\\

These metrics are applied to both the sand and loam continua. 
The resulting spatial distributions of the average saturation error and water table error are shown in Figure~\ref{fig:error}, 
which provide a quantitative assessment of the model discrepancies and complement the qualitative comparisons presented earlier.  

\begin{figure}[H]
\centering
\includegraphics[width=0.9\textwidth]{images/err.png}
\caption{Normalized mean absolute saturation field and water table errors between the Boussinesq–Richards and the 2D Richards models\\
(a) saturation error; (b) water table error. Results are shown for the sand domain over one year.}
\label{fig:error}
\end{figure}


As shown in Figure~\ref{fig:error}, the saturation error remains below \(5\%\) in the sand continuum and \(3\%\) in the loam continuum. 
The larger discrepancy in sand arises from its higher hydraulic conductivity, which allows faster redistribution of water in the unsaturated zone. Since the Boussinesq--Richards approximation neglects part of the vertical fluxes, these effects are more pronounced in sand and lead to slightly larger deviations in the saturation field. In loam, where the conductivity is smaller, water moves more slowly and the neglect of vertical fluxes has a weaker impact, resulting in lower saturation errors.\\

For the table, the error is smaller in sand (below \(6.5\%\)) than in loam (up to \(8\%\)). In sand, the water table position is primarily governed by the overall balance of recharge and discharge, which the Boussinesq-Richards model accurately captures. However, in loam, the lower conductivity amplifies the influence of the capillary zone on water table dynamics. These effects are fully resolved in the Richards model but remain unaccounted for in the Boussinesq-Richards approximation, leading to larger deviations.\\

Taken together, these results underline the conceptual distinction between the Richards and Boussinesq formulations: 
the Richards model accounts explicitly for vertical fluxes and unsaturated dynamics, while the Boussinesq approximation emphasizes lateral flow and simplifies the unsaturated zone. Despite this difference, the BR model reproduces the water table profiles of the 2D Richards model with a maximum deviation of only about \(8\%\), indicating that it provides a numerically consistent and computationally efficient alternative suitable for hydrological applications.\\





%%%%%%%%%%%%%%%%%%%%%%%%%%%%%%%%%%%%%%%%%%%%%%%%%
%%                Outflow                      %%
%%%%%%%%%%%%%%%%%%%%%%%%%%%%%%%%%%%%%%%%%%%%%%%%%

In addition to saturation distributions and water table positions, it is essential to examine the outflow predicted by the two models. 
Outflow represents an integrated hydrological response of the system, as it directly links recharge to discharge and thereby reflects the overall water balance. 
While saturation and water table provide local and spatially distributed insights, outflow is a bulk quantity that is often of primary interest in practical applications such as catchment hydrology and groundwater management. \\

Therefore, comparing the outflow simulated by the Boussinesq–Richards model against that of the reference 2D Richards model constitutes a key step in assessing the reliability of the coupled approach.\\

Figure~\ref{fig:OF} presents the instantaneous outflow rates and cumulative outflows for both sand and loam continua. 

\begin{figure}[H]
\centering
\makebox[\textwidth]{
    \includegraphics[width=0.95\textwidth, height=0.65\textwidth]{images/OF.png}
}
\caption{\centering
Instantaneous and cumulative outflow rates for sand and loam continua, comparing the Boussinesq–Richards and 2D Richards models.\\
(a) and (c) Instantaneous outflow; (b) and (d) Cumulative outflow.
}
\label{fig:OF}
\end{figure}

Figure~\ref{fig:OF} shows close agreement in outflow between the Boussinesq-Richards and Richards models. Both capture similar delays in outflow response to recharge, reflecting unsaturated-zone percolation times. Due to higher hydraulic conductivity, the sand system exhibits faster drainage with sharper, higher outflow peaks. Conversely, the lower-conductivity loam shows delayed peaks and prolonged gradual release, consistent with its stronger retention. Critically, both cases maintain excellent model alignment, confirming the Boussinesq-Richards approach reliably reproduces integrated water-balance dynamics.







\subsubsection{Case 2: test case on a realistic recharge time series.}
\vspace{0.3cm}
%%%%%%%%%%%%%%%%%%%%%%%%%%%%%%%%%%%%%%%%%%%%%%%%%
%%            Realistic test case              %% 
%%%%%%%%%%%%%%%%%%%%%%%%%%%%%%%%%%%%%%%%%%%%%%%%%

In this test case, the simulation domain is defined as a loam slope with a length of 50 meters (corresponding to $x \in [-25, 25]$), a vertical height of 10 meters, and an inclination of $2\%$. To better represent natural surface undulations, the top boundary is defined by a base inclined plane with superimposed sinusoidal perturbations. The recharge dataset used is provided by \textit{Météo-France} (\texttt{DRAINC\_2751\_daily.input})\cite{QuintanaSegui2020SURFEX}.\\

The top boundary of the slope is defined by a base inclined plane with superimposed sinusoidal perturbations to better represent natural surface undulations. Mathematically, the elevation $z_t(x)$ is given by:

\begin{equation}
z_t(x) =
\begin{cases} 
0.2 x + 10 + 0.1 \sin\left(\frac{2 \pi x}{3.5}\right), & 13 \le x \le 77, \\
0.2 x + 10 + 0.05 \sin\left(\frac{2 \pi x}{1.5}\right), & 78 \le x \le 119, \\
0.2 x + 10 + 0.04 \sin\left(\frac{2 \pi x}{0.8}\right), & 120 \le x \le 220, \\
0.2 x + 10, & \text{otherwise}.
\end{cases}
\end{equation}


We first apply the mean recharge value derived from the full dataset, as shown in Figure~\ref{fig:recharge_data} and Figure~\ref{fig:3y_recharge}, to establish a steady-state.\\

All recharge dataset is illustrated in Figure~\ref{fig:recharge_data}.

\begin{figure}[H]
    \centering
    \includegraphics[width=0.6\textwidth]{images/11.png}
    \caption{Visualization of the recharge data}
    \label{fig:recharge_data}
\end{figure}

To reduce computational time, only the first three years of data were used in the simulation. The corresponding recharge time series is shown in Figure~\ref{fig:3y_recharge}.

\begin{figure}[H]
    \centering
    \includegraphics[width=0.6\linewidth]{images/12.png}
    \caption{Recharge from 1957 to 1960}
    \label{fig:3y_recharge}
\end{figure}


%%%%%%%%%%%%%%%%%%%%%%%%%%%%%%%%%%%%%%%%%%%%%%%%%
%%               initial state                 %% 
%%%%%%%%%%%%%%%%%%%%%%%%%%%%%%%%%%%%%%%%%%%%%%%%%


The first three years of recharge data was applied on steady-sate (shown in Figure~\ref{fig:steady_state}.) to initiate the transient simulation.
\vspace{0.5cm}
\begin{figure}[H]
    \centering
    \includegraphics[width=0.8\textwidth]{images/13.jpg}
    \caption{\centering
        Steady-state saturation results.\\
        Top: Boussinesq-Richards;Bottom: Richards.
    }
    \label{fig:steady_state}
\end{figure}


    %%%%%%%%%%%%%%%%%%%%%%%%%%%%%%%%%%%%%%%%%%%%%%%%%
    %%                   video                     %% 
    %%%%%%%%%%%%%%%%%%%%%%%%%%%%%%%%%%%%%%%%%%%%%%%%%

The transient-state simulation is available in the video: 
\href{run:video1.mp4}{\texttt{video1.mp4}}.\\

    %%%%%%%%%%%%%%%%%%%%%%%%%%%%%%%%%%%%%%%%%%%%%%%%%
    %%                Outflow                      %% 
    %%%%%%%%%%%%%%%%%%%%%%%%%%%%%%%%%%%%%%%%%%%%%%%%%
To further analyze the system's dynamic response under realistic conditions, the outflow during the transient phase is evaluated. As in before, both the instantaneous outflow rate and the cumulative outflow volume obtained from the Boussinesq--Richards and 2D Richards models are compared, as shown in Figure~\ref{fig:realistic_outflow}. Similar to Case~1, these two models show excellent agreement when forced by realistic recharge data, accurately replicating the delayed outflow response to recharge events. These results confirm that the Boussinesq–Richards model consistently achieves numerical accuracy while capturing the system's physical dynamics.

\begin{figure}[H]
    \centering
    \makebox[\textwidth]{
    \includegraphics[width=1\textwidth, height=0.35\textwidth]{images/14.png}
    }
    \caption{\centering
        Outflow for the first three years (1957-1960).\\
    }
    \label{fig:realistic_outflow}
\end{figure}

In addition, the mean outflow is 479.978~mm/year for the Boussinesq--Richards model and 467.205~mm/year for the 2D Richards model, corresponding to a difference of only 2.73\%, which indicates very good agreement. The mean outflow is calculated as:
\begin{equation}
    M_{\text{outflow}} = \frac{\text{outflow}}{1 \times L} 
    \times 365 \times 1000
\end{equation}
where $M_{\text{outflow}}$ denotes the mean instantaneous outflow rate (in mm/year), $L$ is the length of the simulation domain ($50$~m), and the unit width of the 2D domain is taken as 1~m.\\


    %%%%%%%%%%%%%%%%%%%%%%%%%%%%%%%%%%%%%%%%%%%%%%%%%
    %%                   error                     %% 
    %%%%%%%%%%%%%%%%%%%%%%%%%%%%%%%%%%%%%%%%%%%%%%%%%
The normalized mean absolute error between the Boussinesq-Richards and 2D Richards models was computed over the simulation period (1957–1960). The saturation error in the loam continuum is shown in Figure~\ref{fig:saterr}.

\begin{figure}[H]
\centering
\includegraphics[width=0.6\textwidth]{images/16.png}
\caption{
Spatial distribution of average saturation error between the Boussinesq-Richards and 2D Richards models (1957–1960).
}
\label{fig:saterr}
\end{figure}
As shown in Figure~\ref{fig:saterr}, the maximum value of this error remains below (\(4\%\)), which confirms the close agreement between the Boussinesq-Richards and 2D Richards models.\\


%%%%%%%%%%%%%%%%%%%%%%%%%%%%%%%%%%%%%%%%%%%%%%%%%
%%                Performance                  %% 
%%%%%%%%%%%%%%%%%%%%%%%%%%%%%%%%%%%%%%%%%%%%%%%%%
\subsection{Performance analysis}


In addition to validating the physical consistency of the Boussinesq–Richards (BR) model with the Richards (R) model, computational efficiency was also assessed. For this purpose, I developed a dedicated script to benchmark execution times of the two models.

\paragraph{Synthetic benchmark cases.} 
The script adopts the same recharge configuration as in Case~1, but with a slightly modified simulation domain: the domain length was reduced from 50~m to 10~m, the slope was increased from 5\% to 10\%, and the height was decreased from 7~m to 5~m. 
This setup ensures comparable hydrological dynamics while allowing for systematic performance evaluation. Table~\ref{tab:runtime} presents the execution times for different grid sizes, highlighting the performance advantage of the BR model over the 2D R model.


\begin{table}[H]
\centering
\caption{Comparison of computational time between the 2D Richards (R) and Boussinesq–Richards (BR) models for different grid sizes. Best BR performance is highlighted in bold.}
\label{tab:runtime}
\begin{tabular}{c|c|c|c}
\hline
Grid size & Model & Time (s) & Speedup (vs R) \\
\hline
\multirow{5}{*}{50 $\times$ 25} 
& R         & 5.237  & -- \\
& BR(1proc) & 1.390 & 3.8$\times$ \\
& BR(2proc) & \textbf{1.266} & 4.1$\times$ \\
& BR(4proc) & 1.372 & 3.8$\times$ \\
& BR(8proc) & 2.461 & 2.1$\times$ \\
\hline
\multirow{5}{*}{100 $\times$ 50} 
& R         & 20.245  & -- \\
& BR(1proc) & 2.981 & 6.8$\times$ \\
& BR(2proc) & 2.187 & 9.3$\times$ \\
& BR(4proc) & \textbf{1.835} & 11.0$\times$ \\
& BR(8proc) & 3.153 & 6.4$\times$ \\
\hline
\multirow{5}{*}{150 $\times$ 75} 
& R         & 48.599  & -- \\
& BR(1proc) & 5.679 & 8.6$\times$ \\
& BR(2proc) & 3.547 & 13.7$\times$ \\
& BR(4proc) & \textbf{2.708} & 17.9$\times$ \\
& BR(8proc) & 3.790 & 12.8$\times$ \\
\hline
\multirow{5}{*}{200 $\times$ 100} 
& R         & 97.370  & -- \\
& BR(1proc) & 9.220 & 10.6$\times$ \\
& BR(2proc) & 5.647 & 17.3$\times$ \\
& BR(4proc) & \textbf{3.921} & 24.8$\times$ \\
& BR(8proc) & 5.188 & 18.8$\times$ \\
\hline
\end{tabular}
\end{table}
 
Table~\ref{tab:runtime} presents the execution times for different grid sizes, highlighting the performance advantage of the BR model over the 2D R model.  
For small grids ($50 \times 25$), the BR model is already about 4 times faster. As the grid resolution increases, the efficiency gap widens dramatically: for $200 \times 100$ cells, the BR model runs nearly 25 times faster than the R model when using four MPI processes.\\

It is worth noting that the efficiency of the BR implementation depends on both grid size and the number of processes.  
For small grids, communication overhead dominates, and fewer processes (or even a single process) yield the best performance.  
In contrast, for larger grids, the computational workload grows significantly, which offsets communication costs and allows more processes to deliver meaningful speedup.  
At the same time, larger grids also imply that both the R and BR models must solve substantially larger systems, making computational efficiency increasingly important.  

\paragraph{Realistic recharge case.}  
To further evaluate performance in a realistic hydrological setting, both models were run using realistic recharge data.  
Based on the considerations above, a grid of $250 \times 50$ was adopted, and the BR model was executed with 4 MPI processes to balance computational workload and communication efficiency.  
The simulations covered two different time spans: a short 3-year period (1957--1960) and the full 61-year record (1957--2018).  
The results, reported in Table~\ref{tab:runtime_realistic}, show that the BR model is more than 18 times faster for the short period, and over 20 times faster for the long-term simulation.  
Notably, for the 61-year case, the R model requires almost 49 minutes, while the BR model completes the same run in less than 2.5 minutes.  
This striking difference demonstrates the practical significance of the BR formulation for long-term hydrogeological studies.  


\begin{table}[H]
\centering
\caption{Execution times for the realistic recharge test case (grid size $250 \times 50$). Results are reported for both the reference 2D Richards (R) model and the Boussinesq–Richards (BR) model with 4 MPI processes.}
\label{tab:runtime_realistic}
\begin{tabular}{c|c|c|c}
\hline
Recharge data & Model & Time & Speedup (vs R) \\
\hline
1957--1960 (3 years) & R  & 143.583 s & -- \\
                     & BR(4proc) & 7.730 s & 18.6$\times$ \\
\hline
1957--2018 (61 years) & R  & 2929.735 s (48.829 min) & -- \\
                      & BR(4proc) & 142.448 s (2.374 min) & 20.6$\times$ \\
\hline
\end{tabular}
\end{table}

\paragraph{Discussion.}
This scaling behavior demonstrates two key aspects:  
\begin{itemize}
    \item \textbf{Numerical stability and robustness:} The BR model avoids convergence issues in many situations where the 2D Richards model struggles, while still reproducing consistent results.  
    \item \textbf{Physical consistency:} The BR formulation preserves the essential dynamics of recharge, discharge, and water table fluctuations, in line with soil hydraulic properties.  
    \item \textbf{Computational efficiency:} Across all tested grid sizes, the BR model consistently outperforms the 2D R model in terms of execution time. Depending on the resolution, the BR simulations are between 4 and 25 times faster than the R simulations (see Table~\ref{tab:runtime}). This substantial gain highlights the advantage of the BR model, which reduces the computational cost while still preserving the key hydrological dynamics.
\end{itemize}

Overall, the BR model provides a numerically stable, physically meaningful, and computationally efficient alternative to the 2D Richards model, making it particularly suitable for large-scale hydrogeological applications.\\



%%%%%%%%%%%%%%%%%%%%%%%%%%%%%%%%%%%%%%%%%%%%%%%%%
%%      Contribution to scientific writing     %%
%%%%%%%%%%%%%%%%%%%%%%%%%%%%%%%%%%%%%%%%%%%%%%%%%
\section{Contribution to scientific writing}

Besides the technical aspects of model development and validation, this internship also provided me with valuable experience in a scientific article writing.  
Producing a structured and rigorous report was an integral part of the project, requiring me to translate numerical results into a clear and convincing narrative. \\ 


My main contributions to the writing process include:
\begin{itemize}
    \item Drafting and structuring methodological sections, covering model testing, validation, and performance analysis;
    \item preparing figures, tables, and captions to support the interpretation of results;
    \item ensuring consistency of notation, equations, and terminology across sections;
    \item articulating physical interpretations of numerical results, highlighting their hydrological relevance;
    \item summarizing validation outcomes and performance insights in a way suitable for potential publication. 
\end{itemize}

This writing effort strengthened my ability to structure scientific arguments, balance technical detail with readability, and make effective use of visual material.\\  

Overall, this part of the internship emphasized the role of writing as a core component of scientific research.  
The skills acquired here will be directly transferable to the preparation of future research articles, my doctoral thesis, and broader contributions to the scientific community.\\





%%%%%%%%%%%%%%%%%%%%%%%%%%%%%%%%%%%%%%%%%%%%%%%%%
%%                Conclusion                   %%
%%%%%%%%%%%%%%%%%%%%%%%%%%%%%%%%%%%%%%%%%%%%%%%%%
\section{Conclusion and perspectives}

This internship has given me valuable insights into how a research project is carried out and advanced within a team environment.
In addition to technical development, I gained experience in interdisciplinary collaboration and teamwork, which are essential skills for pursuing scientific research.
From a mathematical perspective, I found the study of the Boussinesq–Richards model particularly stimulating.
In particular, its strategy of combining vertical and horizontal flow treatments to replace a fully nonlinear Richards system offered me a new perspective on how model complexity can be reduced while still retaining physical realism.
The internship enhanced my ability to programmaticallly generate geometric configurations for simulations, along with debugging, analyzing results, and connecting numerical behavior to physics.\\  

Looking ahead, improving model efficiency while maintaining physical realism is an interest research direction.
The Boussinesq–Richards hybrid model already demonstrates this potential by reducing computational cost without losing hydrological relevance.
Based on this idea, strategies such as model simplification, dimensionality reduction, or hybrid methods that combine detailed process representations with efficient large-scale approximations deserve further exploration.
Another promising direction is extending the MPI implementation of the Boussinesq–Richards model to improve scalability on high-performance computing platforms, which would allow applications to larger and more complex hydrological systems.\\

This internship has thus enriched both my technical and scientific skills, and it reinforced my interest in pursuing further research on numerical modeling, model simplification, and efficient computation for complex physical systems.












\section{Acknowledgements}  
I would like to express my sincere gratitude to my supervisors for their continuous guidance, valuable feedback, and encouragement throughout the internship.  
In particular, I am deeply thankful to Mr. Konstantin Brenner for his technical support in both mathematical aspects and code development, and to Mr. Brown Jean Marçais and Mr. Brown Jean-Raynald de Dreuzy for their insightful advice on the physical interpretation and hydrological theory.  
This work has greatly benefited from their expertise and collaboration.  






\newpage
\bibliographystyle{plain}
\bibliography{biblio}
\end{document}
